\subsubsection{Επίδραση θερμοκρασίας}
\label{subsubsec:temp_effect}
Η πρόσφυση του ελαστικού παρουσιάζει έντονα μη γραμμική εξάρτηση από τη
θερμοκρασία του πέλματος. Σε μια τυπική συμπεριφορά ενός αγωνιστικού ελαστικού
καθώς η θερμοκρασία αυξάνεται, η διαθέσιμη πρόσφυση αρχικά αυξάνεται 
μέχρι ένα βέλτιστο εύρος θερμοκρασίας (παράθυρο λειτουργίας) και στη
συνέχεια ακολουθεί μια πτώση καθώς το ελαστικό εξέρχεται απο το παράθυρο λειτουργίας.
Το παράθυρο λειτουργίας δεν είναι μια καθολική ιδιότητα, 
αλλά εξαρτάται έντονα από τη χημική σύνθεση του ελαστικού το επίπεδο φθοράς κ.α.
Για παράδειγμα, αναφέρεται στην έρευνα των West, Limebeer \cite{West2022}
πως η βέλτιστη πρόσφυση σε ένα αγωνιστικό ελαστικό εμφανίζεται 
σε θερμοκρασία πέλματος περίπου μεταξύ $90\degC$ και
$105\degC$. Επομένως, σε ένα μοντέλο οχήματος που
υπερβαίνει τις ιδεατές συνθήκες η μοντελοποίηση της επίδρασης της θερμοκρασίας
παίζει πολύ σημαντικό ρόλο στις τελικές επιδόσεις του οχήματος. Άλλωστε, σε
αγωνιστικές εφαρμογές, ειδικότερα στο Παγκόσμιο Πρωτάθλημα Formula 1, όπου
τα ελαστικά είναι ιδιαίτερα ευαίσθητα στη θερμοκρασία, μεγάλο μέρος της
στρατηγικής του αγώνα, του σχεδιασμού (βλ. DAS \cite{WikipediaDualAxisSteering}),
των παραμέτρων ευθυγράμμισης κ.α αφορούν τη διαχείριση θερμοκρασίας.

Σε φυσικό επίπεδο, η θερμοκρασία του πέλματος επηρεάζει την πρόσφυση επειδή
μεταβάλλει τις ιξωδοελαστικές ιδιότητες του ελαστομερούς. Η τριβή
ελαστικού και οδοστρώματος σε μοριακό επίπεδο αποδίδεται σε
συνεισφορές που σχετίζονται με τη μοριακή προσκόλληση και με τις απώλειες
παραμόρφωσης του ελαστικού από τις ανωμαλίες του οδοστρώματος (υστέρηση). 
Η μεταβολή της θερμοκρασίας αλλάζει τη δυναμική απόκριση και τους χρόνους 
χαλάρωσης του ελαστομερούς \cite{Grosch1963}. Επιπλέον, η αύξηση της θερμοκρασίας αυξάνει τη
δραστικότητα του πολυμερούς αυξάνοντας τις δυνάμεις προσκόλλησης (chemical adhesion). Η πτώση
της πρόσφυσης με την περαιτέρω αύξηση της θερμοκρασίας αποδίδεται ως επί το πλείστον
στον εκφυλισμό των φυσικών ιδιοτήτων του, ενώ επίσης στο εξωτερικό στρώμα του πέλματος
επιταχύνεται η φθορά και πρόσφυση περιορίζεται από την 
αντοχή του ίδιου του πέλματος.


Είναι πλέον συχνό, τα αγωνιστικά ελαστικά
να μην παρέχονται πλήρως πολυμερισμένα, με αποτέλεσμα στα πρώτα χιλιόμετρα
χρήσης να είναι μαλακότερα έχοντας μεγαλύτερη πρόσφυση (αλλά και ρυθμό φθοράς).
Μετά τον πρώτο κύκλο θερμικής φόρτισης το πέλμα πλέον πολυμερίζεται πλήρως 
(δημιουργία διασταυρούμενων δεσμών -- βουλκανισμός)
και οι ιδιότητές του σταθεροποιούνται παρέχοντας μεγαλύτερη διάρκεια. Σε κάποιες
αγωνιστικές κατηγορίες οι ομάδες επιλέγουν να προετοιμάσουν τα ελαστικά σε δοκιμές (scrubbing in)
ενώ σε άλλες κατηγορίες όπως η Formula 1, χρησιμοποιούνται σχεδόν αποκλειστικά 
στα προκριματικά (qualifying) όπου απαιτούνται οι μέγιστες δυνατές επιδόσεις αλλά
για λίγους μόλις γύρους.
